% Section 3.2, from lectures/L03/README.md: the objectives, the evaluation
% questions, and the next lecture.

\section{Review}
\label{c3:sec:review}

\needspace{6\baselineskip}
\subsection*{What you should be able to do now}
After this chapter, you should:

\begin{itemize}
  \item Be able to create threads statically and dynamically, and choose their priorities.
  \item Understand exactly when a cooperative and a preemptible thread can lose the CPU.
  \item Be able to predict the order in which a set of threads runs.
  \item Understand how timeouts are converted to ticks, and what that does to the length of a sleep.
  \item Understand which context a timer’s expiry function runs in, and what may not be done there.
  \item Be able to defer work from interrupt context to a work queue, and reschedule delayable work.
  \item Understand what \code{native_sim} simulates faithfully about scheduling, and what it does not.
  \item Be able to estimate a thread’s stack needs from static analysis.
\end{itemize}

\needspace{6\baselineskip}
\subsection*{Questions to test yourself}
You should be able to explain:

\begin{enumerate}
  \item When can a cooperative thread lose the CPU, and when can a preemptible one?
  \item Why does \code{k_yield()} sometimes return at once, even though other threads are ready?
  \item Why does \code{k_msleep(30)} return after 40 ms on a 100 Hz tick, while a 500 ms one-shot timer
    fires at exactly 500 ms?
  \item What must not be done in a timer’s expiry function, and why?
  \item Why is the system work queue a cooperative thread, and what does that mean for the work items on it?
  \item When would you use a work queue of your own rather than the system work queue?
  \item Why does \code{while (!flag) \{\}} stall \code{native_sim} forever, when on a microcontroller it
    would only waste power?
  \item Why are Zephyr’s thread analyzer and stack overflow checks not available on \code{native_sim}?
\end{enumerate}

\needspace{6\baselineskip}
\subsection*{Looking ahead}
Chapter~\ref{c4:ch} is Kernel II, communication and application architecture:

\begin{itemize}
  \item Semaphores, mutexes, message queues and events.
  \item Priority inversion, produced and then fixed.
  \item zbus: publish and subscribe between parts of an application.
  \item The State Machine Framework.
\end{itemize}

You will produce priority inversion on purpose, and measure what a mutex does about it.
